%%
%% paper79_temperley_lieb_jones_en.tex
%% Autor: Michael Fuhrmann
%% Projekt: Specialist Mathematics Project
%% Thema: The Temperley-Lieb Algebra and Jones Polynomial: Open Problems
%% Sprache: Englisch
%% Erstellt: 2026-03-12
%% Letzte Änderung: 2026-03-12
%% Build: 164
%%
%% Beschreibung:
%%   Dieses Paper behandelt die Temperley-Lieb-Algebra TL_n(δ), den Jones-
%%   Knoteninvarianten V_L(t), die Khovanov-Homologie als Kategorifizierung,
%%   offene Probleme (Jones-Unknoten-Vermutung, Volumen-Vermutung).
%%   Die Jones-Unknoten-Vermutung und die Volumen-Vermutung sind offen.
%%

\documentclass[12pt,a4paper]{amsart}

%% Pakete für Mathematik, Layout, Encoding und Links
\usepackage{amsmath,amssymb,amsthm,mathtools,enumitem,booktabs,array,hyperref}
\usepackage[utf8]{inputenc}
\usepackage[T1]{fontenc}
\usepackage{geometry}
\geometry{margin=2.5cm}

%% Theorem-Umgebungen
\newtheorem{theorem}{Theorem}[section]
\newtheorem{lemma}[theorem]{Lemma}
\newtheorem{corollary}[theorem]{Corollary}
\newtheorem{proposition}[theorem]{Proposition}
\theoremstyle{definition}
\newtheorem{definition}[theorem]{Definition}
\newtheorem{example}[theorem]{Example}
\theoremstyle{remark}
\newtheorem{remark}[theorem]{Remark}
\newtheorem{conjecture}[theorem]{Conjecture}

%% Mathematische Operatoren und Abkürzungen
\newcommand{\Z}{\mathbb{Z}}
\newcommand{\Q}{\mathbb{Q}}
\newcommand{\C}{\mathbb{C}}
\newcommand{\R}{\mathbb{R}}
\DeclareMathOperator{\tr}{tr}
\DeclareMathOperator{\rank}{rank}
\DeclareMathOperator{\Kh}{Kh}

%% Hyperref-Konfiguration
\hypersetup{
    colorlinks=true,
    linkcolor=blue,
    citecolor=blue,
    urlcolor=blue,
    pdftitle={The Temperley-Lieb Algebra and Jones Polynomial: Open Problems},
    pdfauthor={Michael Fuhrmann}
}

%% -------------------------------------------------------------------------
\begin{document}
%% Abstract VOR \maketitle
%% -------------------------------------------------------------------------
\begin{abstract}
The Temperley--Lieb algebra $\mathrm{TL}_n(\delta)$, introduced by Lieb
and Temperley in the context of statistical mechanics~\cite{TemperleyLieb1971},
provides the algebraic framework for the Jones polynomial $V_L(t)$,
a landmark knot invariant discovered by Vaughan Jones in 1984~\cite{Jones1985}.
We present the diagrammatic calculus of $\mathrm{TL}_n(\delta)$,
the Markov trace and its role in defining $V_L(t)$, the skein relation,
and the main properties of the Jones polynomial.

We then discuss Khovanov homology~\cite{Khovanov2000} as a \emph{categorification}
of the Jones polynomial: a bigraded homology theory whose graded Euler
characteristic recovers $V_L(t)$.

The central open problems are:
\begin{itemize}
\item The \emph{Jones unknot conjecture}: whether $V_L(t) = 1$ implies
      $L$ is the unknot;
\item The \emph{volume conjecture} of Kashaev--Murakami--Murakami:
      whether the exponential growth rate of the coloured Jones polynomial
      equals the hyperbolic volume of the knot complement;
\item Whether Khovanov homology detects the unknot (known to hold by
      Kronheimer--Mrowka~\cite{KronheimerMrowka2011}).
\end{itemize}
\end{abstract}

\title{The Temperley--Lieb Algebra and Jones Polynomial:\\
Open Problems}
\author{Michael Fuhrmann}
\address{Specialist Mathematics Project, Build 165}
\email{specialist-maths@project.local}
\date{2026-03-12}

\maketitle

%% -------------------------------------------------------------------------
\section{Introduction}
%% -------------------------------------------------------------------------

In 1984, while working on von Neumann algebras, Vaughan Jones discovered
a new polynomial invariant of knots and links~\cite{Jones1985}. The Jones
polynomial $V_L(t)$ is a Laurent polynomial in $t^{1/2}$ that is a
topological invariant of the oriented link $L$. Its discovery revolutionised
knot theory, earned Jones the Fields Medal (1990), and opened new connections
to quantum groups, statistical mechanics, conformal field theory, and
3-dimensional quantum topology.

The algebraic underpinning of the Jones polynomial is the
\emph{Temperley--Lieb algebra} $\mathrm{TL}_n(\delta)$, an associative
algebra with generators $\{e_1, \ldots, e_{n-1}\}$ that admits a beautiful
diagrammatic realisation in terms of ``tangle diagrams.'' The Jones
polynomial is constructed via a Markov trace on $\mathrm{TL}_n(\delta)$,
which turns out to be invariant under braid isotopy and Markov moves.

Despite its success in distinguishing many knots, the Jones polynomial
leaves fundamental questions open. Most notably:

\begin{conjecture}[Jones Unknot Conjecture]
\label{conj:Jones-unknot}
If $V_L(t) = 1$, then $L$ is the unknot.
\end{conjecture}

This deceptively simple question — whether the Jones polynomial detects the
unknot — remains unresolved after 40 years.

%% -------------------------------------------------------------------------
\section{The Temperley--Lieb Algebra}
%% -------------------------------------------------------------------------

\subsection{Definition and Relations}

\begin{definition}[Temperley--Lieb Algebra]
Let $\delta \in \mathbb{C}$ (or a commutative ring $R$ with element $\delta$).
The \emph{Temperley--Lieb algebra} $\mathrm{TL}_n(\delta)$ is the unital
associative $\mathbb{C}$-algebra with generators $\{e_1, e_2, \ldots, e_{n-1}\}$
and relations:
\begin{enumerate}[label=\roman*)]
\item $e_i^2 = \delta \cdot e_i$ \hfill (idempotent up to scalar);
\item $e_i e_{i\pm 1} e_i = e_i$ \hfill (Temperley--Lieb relation);
\item $e_i e_j = e_j e_i$ for $|i-j| \geq 2$ \hfill (distant commutativity).
\end{enumerate}
\end{definition}

The algebra $\mathrm{TL}_n(\delta)$ has dimension equal to the Catalan number
$C_n = \frac{1}{n+1}\binom{2n}{n}$.

\begin{theorem}
$\dim_\mathbb{C} \mathrm{TL}_n(\delta) = C_n = \frac{1}{n+1}\binom{2n}{n}$
for generic $\delta$.
\end{theorem}

\begin{proof}
A basis is given by the set of \emph{$(n,n)$-Temperley--Lieb diagrams}:
isotopy classes of planar diagrams consisting of $n$ points on the top,
$n$ points on the bottom, connected by non-crossing arcs within a rectangle.
The number of such diagrams is the Catalan number $C_n$. The multiplication
is given by vertical concatenation (with loops evaluated to $\delta$).
The verification that these diagrams satisfy the relations and are linearly
independent is a combinatorial exercise.
\end{proof}

\subsection{Diagrammatic Calculus}

The Temperley--Lieb algebra has a natural graphical representation:
\begin{itemize}
\item The identity $1 \in \mathrm{TL}_n(\delta)$ is represented by $n$
      vertical strands.
\item The generator $e_i$ is represented by a diagram where strands $i$
      and $i+1$ on the top are connected by a cup, strands $i$ and $i+1$
      on the bottom are connected by a cap, and all other strands are
      vertical.
\item Multiplication corresponds to vertical stacking.
\item A closed loop evaluates to $\delta$.
\end{itemize}

This diagrammatic language makes the relations geometrically transparent:
$e_i^2 = \delta e_i$ corresponds to removing a closed loop, and
$e_i e_{i+1} e_i = e_i$ corresponds to an isotopy of tangles.

\subsection{Inclusion in the Hecke Algebra}

\begin{definition}[Iwahori--Hecke Algebra]
The \emph{Hecke algebra} $H_n(q)$ is the quotient of the group algebra
$\mathbb{C}[B_n]$ (braid group algebra) by the relations
$(T_i - q)(T_i + 1) = 0$, where $T_i$ are the standard braid generators.
\end{definition}

\begin{theorem}[Jones 1987~\cite{Jones1987}]
There is a surjective algebra homomorphism $\rho \colon H_n(q) \to \mathrm{TL}_n(\delta)$
with $\delta = q + q^{-1}$, given by $T_i \mapsto q e_i - 1$.
\end{theorem}

%% -------------------------------------------------------------------------
\section{The Jones Polynomial}
%% -------------------------------------------------------------------------

\subsection{Construction via Markov Trace}

\begin{definition}[Markov Trace]
A \emph{Markov trace} on $\bigcup_n \mathrm{TL}_n(\delta)$ is a family of
linear maps $\tau_n \colon \mathrm{TL}_n(\delta) \to \mathbb{C}$ satisfying:
\begin{enumerate}[label=\roman*)]
\item $\tau_n(ab) = \tau_n(ba)$ (trace property);
\item $\tau_{n+1}(a \cdot e_n) = z \cdot \tau_n(a)$ for all
      $a \in \mathrm{TL}_n(\delta)$ (Markov property),
\end{enumerate}
for some scalar $z$.
\end{definition}

\begin{theorem}[Jones 1985~\cite{Jones1985}]
There exists a unique Markov trace on $\bigcup_n \mathrm{TL}_n(\delta)$
(up to normalisation), given by the ``closure'' of diagrams:
\[
    \tau_n(D) = \delta^{-n} \cdot (\text{number of closed loops in the closure of } D).
\]
\end{theorem}

The Jones polynomial is then defined as follows. Let $L$ be an oriented link
presented as the closure $\hat{\beta}$ of a braid $\beta \in B_n$.
Map $\beta$ to an element $\rho(\beta) \in \mathrm{TL}_n(\delta)$ and take
the Markov trace:

\begin{definition}[Jones Polynomial]
With $\delta = -t^{1/2} - t^{-1/2}$ (so $q = -t^{1/2}$), define
\[
    V_L(t) = \left( -\frac{t^{1/2} + t^{-1/2}}{\delta} \right)^{n-1}
    \cdot \left( -t^{3/4} \right)^{w(\beta)} \cdot \tau_n(\rho(\beta)),
\]
where $w(\beta)$ is the writhe (algebraic length) of $\beta$.
This is independent of the braid representative (by Markov's theorem).
\end{definition}

\subsection{The Skein Relation}

\begin{theorem}[Jones Skein Relation~\cite{Jones1985}]
The Jones polynomial satisfies the skein relation:
\[
    t^{-1} V_{L_+}(t) - t \, V_{L_-}(t) = (t^{1/2} - t^{-1/2}) V_{L_0}(t),
\]
where $L_+$, $L_-$, $L_0$ are three links differing at a single crossing
(positive crossing, negative crossing, and smoothed crossing respectively).
\end{theorem}

\begin{proof}
This follows from the relation $(T_i - q)(T_i + 1) = 0$ in the Hecke algebra,
which gives $T_i - T_i^{-1} = (q - q^{-1}) \cdot 1$. Translated to link
diagrams via the Markov trace, this yields the skein relation.
\end{proof}

\subsection{Properties of the Jones Polynomial}

\begin{theorem}[Basic Properties]
The Jones polynomial $V_L(t)$ satisfies:
\begin{enumerate}[label=\roman*)]
\item $V_{\text{unknot}}(t) = 1$;
\item $V_L(t) \in \Z[t^{1/2}, t^{-1/2}]$ (Laurent polynomial in $t^{1/2}$);
\item $V_{L_1 \sqcup L_2}(t) = (-t^{1/2} - t^{-1/2}) V_{L_1}(t) V_{L_2}(t)$
      (split union);
\item $V_{\bar{L}}(t) = V_L(t^{-1})$ (mirror image);
\item $V_{L_1 \# L_2}(t) = V_{L_1}(t) \cdot V_{L_2}(t)$ (connected sum).
\end{enumerate}
\end{theorem}

\begin{example}[Trefoil Knot]
For the right-handed trefoil $T(2,3)$:
\[
    V_{T(2,3)}(t) = -t^{-4} + t^{-3} + t^{-1}.
\]
For the left-handed trefoil $T(2,-3)$ (mirror):
\[
    V_{T(2,-3)}(t) = -t^{4} + t^{3} + t.
\]
These are different, which proves that the trefoil is chiral (not amphichiral).
\end{example}

%% -------------------------------------------------------------------------
\section{The Jones Unknot Conjecture}
%% -------------------------------------------------------------------------

\begin{conjecture}[Jones Unknot Conjecture~\cite{Jones1985}]
\label{conj:Jones}
If $V_L(t) = 1$, then $L$ is the unknot.
\end{conjecture}

This is the most fundamental open problem about the Jones polynomial.

\begin{remark}
The conjecture is known to hold for:
\begin{itemize}
\item All knots up to 19 crossings (verified computationally).
\item Alternating knots: the Jones polynomial of a non-trivial alternating
      knot cannot be 1 (follows from Tait conjectures, proved by
      Kauffman~\cite{Kauffman1987}, Murasugi~\cite{Murasugi1987},
      Thistlethwaite~\cite{Thistlethwaite1987}).
\item Positive knots (Tanaka~\cite{Tanaka1996}).
\item Fibred knots of certain types (Ghiggini--Hedden, via Heegaard Floer).
\end{itemize}
No counterexample is known.
\end{remark}

\begin{theorem}[Kauffman Bracket~\cite{Kauffman1987}]
\label{thm:Kauffman}
The \emph{Kauffman bracket} $\langle L \rangle \in \Z[A, A^{-1}]$ is defined
by the skein relations:
\[
    \langle \bigcirc \rangle = 1, \quad
    \langle \bigcirc \sqcup D \rangle = (-A^2 - A^{-2})\langle D \rangle,
\]
\[
    \left\langle \vcenter{\hbox{$L_+$}} \right\rangle =
    A \left\langle \vcenter{\hbox{$L_0$}} \right\rangle +
    A^{-1} \left\langle \vcenter{\hbox{$L_\infty$}} \right\rangle.
\]
The Jones polynomial is then
$V_L(t) = (-A)^{-3w(D)} \langle D \rangle \big|_{A = t^{-1/4}}$,
where $w(D)$ is the writhe of a diagram $D$ of $L$.
\end{theorem}

The Kauffman bracket provides an elementary approach to $V_L$, but the
unknot conjecture remains open because we lack a complete characterisation
of when $(-A)^{-3w(D)}\langle D \rangle = 1$.

%% -------------------------------------------------------------------------
\section{Coloured Jones Polynomial}
%% -------------------------------------------------------------------------

\begin{definition}[Coloured Jones Polynomial]
For an irreducible representation $V_N$ of $\mathrm{SU}(2)$ of dimension $N$,
the \emph{$N$-coloured Jones polynomial} $J_N(L; q)$ is the quantum invariant
of $L$ obtained by colouring each component with $V_N$.
For $N=2$, $J_2(L;q)$ is the usual Jones polynomial $V_L(t)$ (up to
normalisation).
\end{definition}

\begin{theorem}[Properties of Coloured Jones]
The coloured Jones polynomials $\{J_N(L;q)\}_{N \geq 1}$ form a
family of polynomial invariants satisfying a linear recurrence relation
(quantum $A$-polynomial), and the sequence completely determines $L$ up
to mirror image (for hyperbolic knots, conjecturally).
\end{theorem}

%% -------------------------------------------------------------------------
\section{The Volume Conjecture}
%% -------------------------------------------------------------------------

\begin{definition}[Kashaev Invariant]
For a knot $K$, the \emph{Kashaev invariant} $\langle K \rangle_N$ at the
$N$-th root of unity $q = e^{2\pi i/N}$ is defined via a specific
$R$-matrix construction~\cite{Kashaev1997}. It was shown to equal the
$N$-coloured Jones polynomial at $q = e^{2\pi i/N}$:
$\langle K \rangle_N = J_N(K; e^{2\pi i/N})$.
\end{definition}

\begin{conjecture}[Volume Conjecture, Kashaev 1997~\cite{Kashaev1997}, Murakami--Murakami 2001~\cite{MurakamiMurakami2001}]
\label{conj:Volume}
For a hyperbolic knot $K$ with hyperbolic volume $\mathrm{Vol}(S^3 \setminus K)$:
\[
    \lim_{N \to \infty} \frac{2\pi \log |J_N(K; e^{2\pi i/N})|}{N}
    = \mathrm{Vol}(S^3 \setminus K).
\]
\end{conjecture}

\begin{example}[Figure-Eight Knot]
For the figure-eight knot $4_1$ (the simplest hyperbolic knot):
\[
    \mathrm{Vol}(S^3 \setminus 4_1) = 3v_3 \approx 2.0298,
\]
where $v_3 \approx 1.0149$ is the volume of the regular ideal tetrahedron.
The volume conjecture has been verified numerically to high precision for
$4_1$ and many other knots.
\end{example}

\begin{theorem}[Murakami--Murakami~\cite{MurakamiMurakami2001}]
For torus knots $T(2,2k+1)$:
\[
    \lim_{N \to \infty} \frac{2\pi \log |J_N(T(2,2k+1); e^{2\pi i/N})|}{N} = 0,
\]
consistent with the volume conjecture (torus knots are not hyperbolic;
their complement has zero simplicial volume).
\end{theorem}

\begin{remark}
The volume conjecture has been proved for torus knots (volume 0), and for
a few specific hyperbolic knots (Ekholm--Ng, 2012 for the figure-eight knot
via a specific ``state sum'' computation), but the general case remains open.
It represents a deep connection between quantum topology and hyperbolic
geometry (Thurston's geometrisation programme).
\end{remark}

%% -------------------------------------------------------------------------
\section{Khovanov Homology}
%% -------------------------------------------------------------------------

\subsection{Categorification of the Jones Polynomial}

\begin{theorem}[Khovanov 2000~\cite{Khovanov2000}]
\label{thm:Khovanov}
There exists a bigraded homology theory $\mathrm{Kh}^{i,j}(L)$, defined
combinatorially from a diagram of $L$, such that:
\begin{enumerate}[label=\roman*)]
\item $\mathrm{Kh}^{i,j}(L)$ is a finitely generated abelian group for
      each $i, j \in \Z$, and is an oriented link invariant;
\item The graded Euler characteristic recovers the Jones polynomial:
      \[
          \sum_{i,j} (-1)^i q^j \dim \mathrm{Kh}^{i,j}(L)
          = (q + q^{-1}) V_L(-q^2);
      \]
\item The Khovanov homology is strictly stronger than the Jones polynomial:
      there exist pairs of links with the same Jones polynomial but
      different Khovanov homology.
\end{enumerate}
\end{theorem}

\subsection{Construction}

The construction of Khovanov homology proceeds as follows:
\begin{enumerate}[label=\arabic*.]
\item Choose a diagram $D$ of $L$ with $n$ crossings; label them $1, \ldots, n$.
\item For each subset $S \subseteq \{1,\ldots,n\}$, form the \emph{$S$-resolution}
      $D_S$ by the 0-smoothing (at crossings not in $S$) and the
      1-smoothing (at crossings in $S$); this gives a collection of circles.
\item Assign to each resolution the vector space $V^{\otimes c(D_S)}$,
      where $V = \Z\{v_+, v_-\}$ is a 2-dimensional module and $c(D_S)$
      is the number of circles.
\item Form the \emph{Khovanov complex} $C^*(D)$ with differential
      $d \colon C^i \to C^{i+1}$ given by multiplication and
      comultiplication maps (from the Frobenius algebra structure on $V$).
\item $\mathrm{Kh}^*(D) = H^*(C^*(D))$ is the Khovanov homology.
\end{enumerate}

\begin{theorem}[Bar-Natan 2002~\cite{BarNatan2002}]
The Khovanov chain complex $C^*(D)$ is well-defined up to chain homotopy
equivalence as a complex of abelian groups (independent of diagram choices).
\end{theorem}

\subsection{Khovanov Homology Detects the Unknot}

\begin{theorem}[Kronheimer--Mrowka 2011~\cite{KronheimerMrowka2011}]
\label{thm:KM}
Khovanov homology detects the unknot: if $\mathrm{Kh}(L) \cong \mathrm{Kh}(\text{unknot})$,
then $L$ is the unknot.
\end{theorem}

\begin{proof}[Proof strategy]
Kronheimer and Mrowka use instanton Floer homology. They define a variant
$\mathrm{Kh}'(L)$ (``reduced Khovanov homology with $\Z/2$-coefficients'')
and show it is related to instanton knot homology $I^\#(L)$ via a spectral
sequence. The instanton homology detects the unknot (since $\dim I^\#(K) = 1$
iff $K$ is the unknot), and the spectral sequence from $\mathrm{Kh}'$ to $I^\#$
collapses at the $E_2$ page for the unknot.
\end{proof}

Theorem~\ref{thm:KM} shows that Khovanov homology is strictly stronger
than the Jones polynomial. The following remains open:

\begin{conjecture}[Jones Polynomial Detects Unknot, Conjecture~\ref{conj:Jones}]
If $V_L(t) = 1$, then $L$ is the unknot.
\end{conjecture}

%% -------------------------------------------------------------------------
\section{Odd Khovanov Homology}
%% -------------------------------------------------------------------------

\begin{definition}[Odd Khovanov Homology~\cite{OzsvathRasmussenSzabo2013}]
The \emph{odd Khovanov homology} $\mathrm{Kh}_{\mathrm{odd}}(L)$ is a variant
of Khovanov homology defined using a different Frobenius algebra structure
(with an exterior algebra instead of the symmetric algebra on $V$).
It satisfies:
\begin{enumerate}[label=\roman*)]
\item $\mathrm{Kh}_{\mathrm{odd}}^{i,j}(L)$ is a link invariant;
\item Over $\Q$: $\mathrm{Kh}_{\mathrm{odd}}(L; \Q) \cong \mathrm{Kh}(L; \Q)$
      (same rational homology);
\item Over $\Z$: $\mathrm{Kh}_{\mathrm{odd}}(L)$ and $\mathrm{Kh}(L)$
      can differ (different $\Z/2$-torsion).
\end{enumerate}
\end{definition}

The relationship between $\mathrm{Kh}(L)$ and $\mathrm{Kh}_{\mathrm{odd}}(L)$
is governed by a long exact sequence:

\begin{theorem}[Ozsváth--Rasmussen--Szabó~\cite{OzsvathRasmussenSzabo2013}]
There is a natural transformation $\mathrm{Kh} \to \mathrm{Kh}_{\mathrm{odd}}$
fitting into a long exact sequence relating the two theories with
$\Z/2$-coefficients.
\end{theorem}

%% -------------------------------------------------------------------------
\section{The Rasmussen \texorpdfstring{$s$}{s}-Invariant from Khovanov Homology}
%% -------------------------------------------------------------------------

\begin{theorem}[Rasmussen 2010~\cite{Rasmussen2010}]
From the Khovanov complex one can extract a concordance invariant
$s(L) \in 2\Z$, the \emph{Rasmussen $s$-invariant}, with properties:
\begin{enumerate}[label=\roman*)]
\item $s$ is a homomorphism on the concordance group;
\item $|s(K)|/2 \leq g_4(K)$;
\item $s$ gives a purely combinatorial proof that torus knots
      $T(p,q)$ have $g_4 = (p-1)(q-1)/2$, previously proved only
      via gauge theory (Kronheimer--Mrowka).
\end{enumerate}
\end{theorem}

This is one of the most striking applications of Khovanov homology:
it gives a combinatorial proof of a result that was previously accessible
only through the deep machinery of gauge theory.

%% -------------------------------------------------------------------------
\section{Open Problems}
%% -------------------------------------------------------------------------

\begin{conjecture}[Jones Unknot Conjecture, Conjecture~\ref{conj:Jones}]
$V_L(t) = 1$ if and only if $L$ is the unknot.
\end{conjecture}

\begin{conjecture}[Volume Conjecture, Conjecture~\ref{conj:Volume}]
For a hyperbolic knot $K$:
\[
    \lim_{N \to \infty} \frac{2\pi \log |J_N(K; e^{2\pi i/N})|}{N}
    = \mathrm{Vol}(S^3 \setminus K).
\]
\end{conjecture}

\begin{conjecture}[Coloured Volume Conjecture]
\label{conj:col-volume}
For the $N$-coloured Jones polynomial evaluated at $e^{2\pi i /N}$,
the phase also encodes the Chern--Simons invariant:
\[
    J_N(K; e^{2\pi i /N}) \sim C \cdot N^{3/2} \cdot \exp\left(
    \frac{N}{2\pi}\bigl(\mathrm{Vol}(S^3 \setminus K)
    + i \cdot \mathrm{CS}(S^3 \setminus K)\bigr) \right),
\]
where $\mathrm{CS}$ is the Chern--Simons invariant.
\end{conjecture}

\begin{conjecture}[AJ Conjecture, Garoufalidis~\cite{Garoufalidis2004}]
\label{conj:AJ}
The quantum $A$-polynomial (the recurrence relation for the coloured Jones
polynomial $\{J_N(K;q)\}$) agrees with the classical $A$-polynomial of
$K$ (defined from the $\mathrm{SL}(2,\C)$ representation variety of
$\pi_1(S^3 \setminus K)$) after specialisation $q \to 1$.
\end{conjecture}

\begin{conjecture}[Functoriality of Khovanov Homology]
\label{conj:Kh-func}
Khovanov homology can be made into a fully extended topological quantum
field theory (TQFT) in the sense of Lurie.
\end{conjecture}

\begin{conjecture}[Odd Khovanov vs.\ Even Khovanov]
There exist prime knots $K$ for which $\mathrm{Kh}(K; \Z)$ and
$\mathrm{Kh}_{\mathrm{odd}}(K; \Z)$ have different ranks over $\Z$
(not just different torsion).
\end{conjecture}

\begin{conjecture}[Knot Floer Homology vs.\ Khovanov Homology]
The spectral sequence from $\mathrm{Kh}(K)$ to $\widehat{HFK}(K)$
(conjectured by Dunfield--Gukov--Rasmussen, related to doubly-graded HOMFLY
homology) exists and yields structural relationships between symplectic
and combinatorial knot invariants.
\end{conjecture}

%% -------------------------------------------------------------------------
\section{Conclusion}
%% -------------------------------------------------------------------------

The Temperley--Lieb algebra and the Jones polynomial represent a confluence
of quantum algebra, diagrammatic topology, and low-dimensional manifold theory.
The Jones polynomial distinguishes the trefoil from its mirror image, detects
many knots that Alexander polynomial cannot, and spawned an industry of
new invariants (HOMFLY, Kauffman bracket, quantum groups, topological quantum
field theories). Yet the simplest question — does $V_L = 1$ imply $L$ is
the unknot? — remains completely open.

Khovanov homology, as a categorification of the Jones polynomial, is
strictly more powerful (it detects the unknot, by Kronheimer--Mrowka),
and its structure organises an enormous amount of information about knot
concordance (via the $s$-invariant). The volume conjecture connects
quantum invariants to hyperbolic geometry in a profound way whose full
mathematical content is still being explored.

\begin{thebibliography}{99}

\bibitem{TemperleyLieb1971}
H.~N.~V.~Temperley, E.~H.~Lieb,
\textit{Relations between the `percolation' and `colouring' problem and
other graph-theoretical problems associated with regular planar lattices:
some exact results for the `percolation' problem},
Proc.\ Roy.\ Soc.\ London Ser.\ A \textbf{322} (1971), 251--280.

\bibitem{Jones1985}
V.~F.~R.~Jones,
\textit{A polynomial invariant for knots via von Neumann algebras},
Bull.\ Amer.\ Math.\ Soc.\ \textbf{12} (1985), 103--111.

\bibitem{Jones1987}
V.~F.~R.~Jones,
\textit{Hecke algebra representations of braid groups and link polynomials},
Ann.\ of Math.\ \textbf{126} (1987), 335--388.

\bibitem{Kauffman1987}
L.~H.~Kauffman,
\textit{State models and the Jones polynomial},
Topology \textbf{26} (1987), 395--407.

\bibitem{Murasugi1987}
K.~Murasugi,
\textit{Jones polynomials and classical conjectures in knot theory},
Topology \textbf{26} (1987), 187--194.

\bibitem{Thistlethwaite1987}
M.~B.~Thistlethwaite,
\textit{A spanning tree expansion of the Jones polynomial},
Topology \textbf{26} (1987), 297--309.

\bibitem{Khovanov2000}
M.~Khovanov,
\textit{A categorification of the Jones polynomial},
Duke Math.\ J.\ \textbf{101} (2000), 359--426.

\bibitem{BarNatan2002}
D.~Bar-Natan,
\textit{On Khovanov's categorification of the Jones polynomial},
Algebr.\ Geom.\ Topol.\ \textbf{2} (2002), 337--370.

\bibitem{Rasmussen2010}
J.~Rasmussen,
\textit{Khovanov homology and the slice genus},
Invent.\ Math.\ \textbf{182} (2010), 419--447.

\bibitem{KronheimerMrowka2011}
P.~B.~Kronheimer, T.~S.~Mrowka,
\textit{Khovanov homology is an unknot-detector},
Publ.\ Math.\ Inst.\ Hautes Études Sci.\ \textbf{113} (2011), 97--208.

\bibitem{Kashaev1997}
R.~M.~Kashaev,
\textit{The hyperbolic volume of knots from the quantum dilogarithm},
Lett.\ Math.\ Phys.\ \textbf{39} (1997), 269--275.

\bibitem{MurakamiMurakami2001}
H.~Murakami, J.~Murakami,
\textit{The colored Jones polynomials and the simplicial volume of a knot},
Acta Math.\ \textbf{186} (2001), 85--104.

\bibitem{OzsvathRasmussenSzabo2013}
P.~Ozsváth, J.~Rasmussen, Z.~Szabó,
\textit{Odd Khovanov homology},
Algebr.\ Geom.\ Topol.\ \textbf{13} (2013), 1465--1488.

\bibitem{Garoufalidis2004}
S.~Garoufalidis,
\textit{On the characteristic and deformation varieties of a knot},
Geom.\ Topol.\ Monogr.\ \textbf{7} (2004), 291--309.

\bibitem{Tanaka1996}
T.~Tanaka,
\textit{Maximal Thurston--Bennequin numbers of 2-bridge links},
Topology Appl.\ \textbf{68} (1996), 1--15.

\end{thebibliography}

\end{document}
